\section{Proposed Benchmark: FinIF}

FinIF evaluates whether an LLM can turn supplied financial materials into an instruction-compliant work product.
Each item consists of a context pack, a workflow-specific query, a set of instruction-following constraints, and an evaluation protocol.
The response is assessed along two separate axes: whether it follows the requested constraints, and whether it is a useful financial answer.
The first axis is the main target of FinIF; the second is kept as an auxiliary diagnostic.

\begin{table}[t]
\centering
\small
\resizebox{\linewidth}{!}{
\begin{tabular}{lcccl}
\toprule
Benchmark & Domain & Test & Train & Main evaluation \\
\midrule
IFEval \citep{zhou2023ifeval} & General & 541 & -- & Rule-based IF checks \\
AgentIF \citep{qi2025agentif} & Agentic & 707 & -- & Code, LLM, hybrid checks \\
SciIF \citep{su2026sciif} & Science & 334 & 910 & Correctness + constraint compliance \\
FIFE \citep{matlin2025fife} & Finance & 88 & -- & Chainable financial verifiers \\
FinIF & Finance & 300 & 1,064 & Rule + LLM audit for workflows \\
\bottomrule
\end{tabular}
}
\caption{Comparison with representative instruction-following benchmarks. FinIF focuses on financial workflow deliverables and provides both a hard test set and training data.}
\label{tab:benchmark_compare}
\end{table}

\subsection{Task Taxonomy}

FinIF uses five workflow stages that cover a broad finance task lifecycle: Intake and Profiling; Research and Due Diligence; Decision and Structuring; Risk and Compliance Review; and Execution, Monitoring, Reporting, and Operations.
These workflows are decomposed into 38 task types.
Examples include KYC onboarding, loan application intake, financial statement analysis, credit due diligence, investment memo drafting, DCF valuation, suitability review, AML red-flag review, reconciliation, board or regulatory reporting, and remediation tracking.
This workflow-first design keeps each prompt tied to a plausible financial deliverable rather than an isolated instruction template.

\begin{table}[t]
\centering
\small
\resizebox{\linewidth}{!}{
\begin{tabular}{lrr}
\toprule
Workflow & Train & Test \\
\midrule
Intake and Profiling & 168 & 54 \\
Research and Due Diligence & 196 & 45 \\
Decision and Structuring & 224 & 63 \\
Risk and Compliance Review & 224 & 64 \\
Execution, Monitoring, Reporting, and Operations & 252 & 74 \\
\midrule
Total & 1,064 & 300 \\
\bottomrule
\end{tabular}
}
\caption{Workflow coverage of FinIF-Train and FinIF-Test.}
\label{tab:workflow}
\end{table}

\subsection{Data Construction}

The FinIF construction pipeline follows a workflow-first sequence:
\begin{quote}
\small
workflow $\rightarrow$ task $\rightarrow$ work product $\rightarrow$ context $\rightarrow$ query $\rightarrow$ constraints $\rightarrow$ evaluation.
\end{quote}

\paragraph{Context acquisition.}
We construct context packs from public primary sources, synthetic private records, or mixed packs.
Public sources include regulatory guidance, public filings, financial statements, and similar primary materials.
Private customer, borrower, account, and internal-policy records are synthetic to avoid exposing sensitive data while preserving realistic workflow triggers.

\paragraph{Query synthesis.}
Each item starts from a target work product, such as a KYC review note, credit memo, investment committee memo, suitability review note, compliance escalation report, reconciliation exception report, or board report.
The query asks the model to produce this deliverable under explicit requirements about format, evidence, quantitative reasoning, decision boundaries, and missing-information handling.

\paragraph{Constraint annotation.}
FinIF evaluates explicit, checkable requirements rather than implicit best practices.
Each constraint receives a fine-grained tag, a family label, an evaluation axis, and a checker route.
The paper-facing taxonomy contains five families: Format and Presentation, Evidence and Grounding, Decision and Boundary, Quantitative Verification, and Required Content Coverage.
For evaluation, content-coverage constraints that primarily diagnose answer completeness are kept separate from instruction-following constraints when computing the main IF metrics.

\paragraph{Evaluation generation.}
For each constraint, we generate an evaluation protocol that routes mechanically checkable requirements to rule checkers and semantic requirements to an LLM judge.
Rule-checked constraints include word ranges, decimal places, table columns, required labels, first-line formats, and ordered-list counts.
LLM-judged constraints include source-label placement, fact-inference separation, evidence-rule-action chains, and boundary reasoning.

\paragraph{Quality control.}
Before inclusion, each item is validated for schema consistency, train/test separation, context availability, checker coverage, and parseable evaluator output.
The reported FinIF-Test runs reach 100\% decision coverage, meaning every instruction-following constraint is assigned a pass/fail decision.

\subsection{Data Statistics}

FinIF-Train contains 1,064 workflow-grounded examples across all 38 task types, with 28 examples per task and 7,851 annotated constraints.
FinIF-Test contains 300 hard items, 3,134 instruction-following constraints, and 1,904 additional diagnostic constraints used for coverage, quality, and finance-validity analysis.
The test set has an average of 10.45 instruction-following constraints per item, with a minimum of 7 and a maximum of 16.
Among the 3,134 test IF constraints, 1,346 are rule-checked and 1,788 are LLM-judged.

\begin{table}[t]
\centering
\small
\resizebox{\linewidth}{!}{
\begin{tabular}{lrr}
\toprule
Constraint family & Test IF & Share \\
\midrule
Format and Presentation & 1,391 & 44.38\% \\
Evidence and Grounding & 1,085 & 34.62\% \\
Decision and Boundary & 280 & 8.93\% \\
Quantitative Verification & 377 & 12.03\% \\
Required Content Coverage & 1 & 0.03\% \\
\midrule
Total & 3,134 & 100.00\% \\
\bottomrule
\end{tabular}
}
\caption{Instruction-following constraint distribution in FinIF-Test. Content coverage is mainly used as an auxiliary diagnostic rather than the primary IF target.}
\label{tab:constraints}
\end{table}
